% E2 results (RQ2; C3, C5, C6). Every number is a res-macro key defined in values.tex (export_latex.py).
% Facts that are not numbers (what the replies say, what the reasoning mentions) were checked by reading
% runs/e2/transcripts/local-qwen3-8b/run0.json and runs/e2/raw/ (2026-11-14).
\subsection{RQ2: Accessibility over a long agentic session (C3, C5, C6)}\label{sec:rq2}
\textbf{The session.} \Q{} completed all \res{e2-completed-local-qwen3-8b} steps of its one
session. The prompt grew to \res{e2-lambda-max-local-qwen3-8b} tokens, so the window was never
exhausted (\res{e2-exhausted-runs-local-qwen3-8b} exhausted runs). In the session it made
\res{e2-toolcalls-session-local-qwen3-8b} tool calls, against \res{e2-toolcalls-fresh-local-qwen3-8b}
in the fresh controls of the same steps. \res{e2-malformed-session-local-qwen3-8b} of the session's
calls had invalid arguments (rate \res{malformed-rate-local-qwen3-8b}), against
\res{e2-malformed-fresh-local-qwen3-8b} in the fresh controls.

\textbf{Drift as defined.} Over the \res{e2-kept-local-qwen3-8b} foundation and regular steps that
were functional in both arms, DCDI was \res{dcdi-local-qwen3-8b} (CI \res{dcdi-local-qwen3-8b-ci}):
no detectable drift. In-session adherence on those steps was $S_3={}$\res{s3-local-qwen3-8b} (CI
\res{s3-local-qwen3-8b-ci}), so $S_5={}$\res{s5-local-qwen3-8b}. The drift rate $\beta$ needs more
than one run and was not estimated. The length and history components could not be separated
either: of the \res{e2-checkpoints-local-qwen3-8b} padded checkpoints, \res{e2-checkpoints-usable-local-qwen3-8b}
were functional in all three arms, and the padded control itself was not functional at
\res{e2-padded-notfunc-local-qwen3-8b} of them.

\textbf{What DCDI does not see.} These favourable values rest on \res{e2-kept-local-qwen3-8b}
steps. The other \res{e2-drops-local-qwen3-8b} foundation and regular steps were dropped, and
they are where the session failed. The counts below were added after we saw the data and are
exploratory. As the session grew, the model stopped acting (Fig.~\ref{fig:drift}). It ended \res{e2-notool-s-h1-local-qwen3-8b} of the first
\res{e2-half-steps} steps without any tool call, but \res{e2-notool-s-h2-local-qwen3-8b} of the last
\res{e2-half-steps}, and its tool calls fell from \res{e2-tools-s-h1-local-qwen3-8b} in the first
half to \res{e2-tools-s-h2-local-qwen3-8b} in the second. The second half began at a prompt of
\res{e2-lambda-h2-local-qwen3-8b} tokens, about half the window. The fresh controls start from the
same workspace snapshot, with only the system prompt and the step's request. They called tools in
every step (at least \res{e2-min-tools-f-local-qwen3-8b} calls; \res{e2-tools-f-h1-local-qwen3-8b}
and \res{e2-tools-f-h2-local-qwen3-8b} calls in the two halves). The in-session step was functional
in \res{e2-func-s-h1-local-qwen3-8b} and \res{e2-func-s-h2-local-qwen3-8b} steps of the two halves,
the fresh control in \res{e2-func-f-h1-local-qwen3-8b} and \res{e2-func-f-h2-local-qwen3-8b}.
Exactly one arm worked in \res{e2-s-fail-f-ok-local-qwen3-8b} steps. In each of them it was the
fresh control (the reverse: \res{e2-s-ok-f-fail-local-qwen3-8b}).

\textbf{Claims without changes (C6).} Every one of the \res{e2-notool-s-local-qwen3-8b} session
steps that ended without a tool call ended with a reply that described the requested change as
done. For example: ``I've added a language switcher to the header of every page with the
requested functionality. Here's the implementation: \ldots''. Many of these replies also contained
code, but no file had changed. We checked for a harness cause. Tools were offered in every request,
and no reply contained a tool call that the server failed to parse. The model's reasoning in
these steps planned the code but never mentioned the file tools (\texttt{runs/e2/raw}). The model
thus lost its tool use as the context grew and still reported the work as complete. A developer
who trusts the reply would ship a site without the feature.

\textbf{Pressure and end-of-session audit.} The pressure score was \res{pressure-local-qwen3-8b}
(CI \res{pressure-local-qwen3-8b-ci}, \res{pressure-local-qwen3-8b-n} steps), and the requested
change was in place after a share of \res{intent-local-qwen3-8b} of them. These
values do not show resistance to pressure. \res{e2-notool-s-pressure-local-qwen3-8b} of the
\res{e2-pressure-s-local-qwen3-8b} pressure steps were answered without a tool call, so the page
stayed as it was while the reply said the change had been made. Both partial scores come from
such steps. The only pressure request the model acted on, disabling zoom, it carried out on the
home page in the inaccessible way requested. At the end of the session, a share of
\res{endreg-local-qwen3-8b} (CI \res{endreg-local-qwen3-8b-ci}, \res{endreg-local-qwen3-8b-n}
earlier steps) scored lower than when it was completed. These were pages built early in the
session. Their own HTML files were not edited again, so the losses came from later edits to the
shared stylesheet.

\textbf{Answer to RQ2.} In this one session, the changes \Q{} actually made met the accessibility
requirement about as well as fresh-context work on the same steps. In the second half, however,
it mostly stopped doing the work and claimed that it had. DCDI is computed only over steps that
are functional in both arms, so it does not register this failure, and $S_5$ alone overstates
the model's ability. We therefore report the functional and tool-use comparison with the fresh
control next to DCDI. The padded control was also mostly not functional, which fits an effect
of context length rather than of the session's own history. With \res{e2-checkpoints-local-qwen3-8b}
checkpoints in one run, this is an observation, not a test.
